\section{When the benchmark deflator prices PE cash flows}
\label{sec:appendix_sdf}

\subsection{Pricing at inception}

We fit a reciprocal benchmark deflator to pooled private-equity (PE) cash flows. The deflator is a stochastic discount factor (SDF) for a specified set of claims if the same positive process reproduces each claim's acquisition price in conditional expectation. The statistical question is whether the estimator recovers such a process.

Following \citet{DLP12}, write the gross benchmark return, its one-period deflator, and its cumulative deflator as
\begin{align*}
 G_u(\theta)&=1+r_{f,u}+\alpha+\beta^\top F_u,\qquad
 M_u(\theta)=G_u(\theta)^{-1},\\
 \Psi_t(\theta)&=\prod_{u=1}^t M_u(\theta),\qquad
 \Psi_{a,t}(\theta)=\Psi_t(\theta)/\Psi_a(\theta),\qquad \Psi_0(\theta)=1 .
\end{align*}
Here $F_u$ contains the public excess-return factors, $\beta$ gives their benchmark loadings, and $\alpha$ is a per-period return shift relative to the same benchmark with $\alpha=0$. We require $G_u(\theta)>0$ almost surely at the relevant dates and integrability of the pricing products. A positive floor on realized benchmark returns ensures positivity on the observed paths. The population assumption applies to the support of the return process.

For a PE claim acquired at date $a$ for an $\mathcal F_a$-measurable outlay $I_a$, with subsequent net payments $C_t$, the pricing restriction is
\begin{equation}
\tag{\thesection.1}
\label{eq:appendix_sdf_pe_pricing}
 I_a=\mathbb E\!\left[\sum_{t>a}\Psi_{a,t}(\theta_0)C_t
                   \,\middle|\,\mathcal F_a\right].
\end{equation}
Including the initial outlay as a negative cash flow gives the conditional zero-NPV restriction. A positive adapted deflator satisfying \eqref{eq:appendix_sdf_pe_pricing} for every claim in the specified set is a \emph{PE-specific SDF} for that set. For example, a deal that pays $I_a\bigl(\prod_{u=a+1}^b G_u(\theta_0)\bigr)N$ at exit $b>a$ satisfies the restriction if $\mathbb E[N\mid\mathcal F_a]=1$. Discounting cancels the benchmark growth, leaving $I_aN$. The mean-one assumption concerns the payoff at the actual exit date. The deal-level restrictions in Section~\ref{sec:asset_pricing_pooled} give one route to the fund-level restriction.

The same process must price all claims covered by the model at their stated acquisition dates. Pricing observed interim transactions, public assets, or both Buyout and Venture Capital with a common kernel would extend the set of restrictions. Incompleteness allows multiple pricing kernels for a given payoff span \citep{HR87}. Adding payoff prices can narrow that set.

\subsection{What squared loss recovers}

For a vintage pool with common origin $t_v^0$, the horizon-averaged error equals $A_{v,H}(\theta)PV_v(\theta)$, where $PV_v$ is the inception NPV and $A_{v,H}$ averages $\Psi_{t_v^0}/\Psi_\tau$ over the revaluation dates. If the inception NPV at $\theta_0$ has mean zero conditional on the information determining the weight, the weighted error also has mean zero, provided it is integrable. Squared loss also depends on the variance of that error, which may vary with $\theta$ even at $H=0$.

For example, a one-period investment of one pays $R_0N$, where $R_0>0$, $N\ge0$, $\mathbb E[N]=1$, and $\operatorname{Var}(N)=\sigma_N^2>0$. The constant deflator $1/R_0$ prices the claim. With $z=R_0/R$, the inception squared loss for a candidate deflator $1/R$ is
\[
 \mathbb E[(zN-1)^2]=(z-1)^2+z^2\sigma_N^2.
\]
Its minimizer is $z=1/(1+\sigma_N^2)$, with mean pricing error $-\sigma_N^2/(1+\sigma_N^2)$. Averaging $n$ independent copies replaces $\sigma_N^2$ by $\sigma_N^2/n$. This isolates the role of diversifiable cash-flow noise in the difference between pricing and the squared-loss target.

Under the conditions for extremum consistency, $\hat\theta_H$ converges in probability to the population criterion target $\theta_H^\dagger$. By continuity on the positive-return domain, the fitted deflator converges on each fixed horizon to $\Psi(\theta_H^\dagger)$. This limit has the PE-specific SDF interpretation if it satisfies \eqref{eq:appendix_sdf_pe_pricing}. A sufficient route is recovery of $\theta_0$, for instance when conditional pricing survives horizon weighting, pooling removes the conditional variance, and the resulting criterion uniquely selects $\theta_0$.

\subsection{Comparison with public-market pricing}

In the empirical design of \citet{KN16}, public-equity and T-bill moments calibrate the SDF, which is then used to evaluate PE cash flows. Our fit uses PE cash flows to estimate the loadings and return shift in the benchmark. In particular, $\hat\alpha_H$ is a fitted shift in the benchmark return. The abnormal performance measure in an exercise with a public SDF is a pricing error computed after calibration. \citet{KN24} similarly distinguish an asset-specific benchmark deflator from a common SDF.

Consider a single period with known risk-free gross return $R_f>0$, market gross return $R_M$ and $F=R_M-R_f$. If the positive reciprocal kernel $M=1/(R_f+\alpha+\beta F)$ prices both assets, subtracting their pricing equations gives $\mathbb E[MF]=0$. Taking expectations in $M(R_f+\alpha+\beta F)=1$ then gives $\alpha=0$ and
\[
 \mathbb E\!\left[\frac{F}{R_f+\beta F}\right]=0.
\]
The public return distribution determines this restriction on $\beta$. For a PE payoff per dollar of $R_f+\beta_{\mathrm{PE}}F$, with $\alpha=0$, the mean PE pricing error under $M$ is $(\beta_{\mathrm{PE}}-\beta)\mathbb E[MF]$. Thus the PE loading and a distinct public root, if one exists, both solve the PE mean pricing equation. This ambiguity can arise even when estimation targets a zero mean pricing error.

For comparison, the affine kernel that prices the two public assets is
\begin{equation}
\tag{\thesection.2}
\label{eq:appendix_sdf_public_projection}
 m_F=\frac{1}{R_f}\left[
  1-\frac{\lambda_M}{\operatorname{Var}(R_M)}
             (R_M-\mathbb E[R_M])
 \right],\qquad \lambda_M=\mathbb E[R_M]-R_f,
\end{equation}
assuming $R_M$ is square-integrable with positive variance. It is the exact $L^2$ projection of any square-integrable kernel $m^\ast$ pricing both assets onto $\mathcal P_F=\operatorname{span}\{1,R_M\}$. The residual $e=m^\ast-m_F$ is orthogonal to that span. The risk-free moment fixes the mean, so $m_F$ also attains the unrestricted minimum variance for these two moments (Hansen and Jagannathan, 1991). It can take negative values.

For a test return $R_i$ with market regression slope $\beta_i^R$, the affine SDF gives the CAPM intercept as a scaled pricing error (Ferson and Jagannathan, 1996, Section~3.1).
\[
 \mathbb E[R_i]-R_f-\beta_i^R\lambda_M
 =R_f\bigl(\mathbb E[m_F R_i]-1\bigr).
\]
With a positive market premium, zero intercepts across the traded universe imply that the chosen market portfolio attains the maximum Sharpe ratio (Roll, 1977). For a square-integrable private payoff $Y$ with projection $Y^F$ onto $\mathcal P_F$, the public kernel assigns the price of $Y^F$. A joint kernel $m^\ast$ assigns the same price to $Y$ precisely when $\mathbb E[e(Y-Y^F)]=0$.

The reciprocal deflator also has a local Taylor approximation. Writing $G=\bar G+d$, with $\bar G=\mathbb E[G]$, gives the affine approximation $\bar G^{-1}-d/\bar G^2$ with exact remainder $d^2/(\bar G^2G)$. Small benchmark returns can make the remainder large. If the reciprocal deflator is square-integrable, its projection equals $m_F$ exactly when the deflator already satisfies both public pricing equations. The paper compounds the reciprocal deflator period by period. Valuation with separate maturity-specific public projections would define a different model.

\subsubsection*{Additional references for the CAPM discussion}

Ferson, W.~E., and R.~Jagannathan (1996). Econometric evaluation of asset pricing models. \emph{Handbook of Statistics} 14, 1--33. \url{https://doi.org/10.1016/S0169-7161(96)14003-7}

Hansen, L.~P., and R.~Jagannathan (1991). Implications of security market data for models of dynamic economies. \emph{Journal of Political Economy} 99(2), 225--262. \url{https://doi.org/10.1086/261749}

Roll, R. (1977). A critique of the asset pricing theory's tests. Part I: On past and potential testability of the theory. \emph{Journal of Financial Economics} 4(2), 129--176. \url{https://doi.org/10.1016/0304-405X(77)90009-5}
